\documentclass[10pt]{article}

\usepackage[utf8]{inputenc}

\usepackage[T1]{fontenc}

\usepackage{amsmath}

\usepackage{amsfonts}

\usepackage{amssymb}

\usepackage[version=4]{mhchem}

\usepackage{stmaryrd}

\usepackage{bbold}

\usepackage{mathrsfs}

\usepackage{graphicx}

\usepackage[export]{adjustbox}

\graphicspath{ {./images/} }



\begin{document}

о пределяется решением $\left\{e_{\alpha}(t)\right\}$ системы



\[

d u_{\alpha}=\left(\omega_{\alpha}^{\beta}\right)_{x(t)}(\dot{x}(t)) u_{\beta}

\]



при начальных условиях $u_{\alpha}(0)=e_{\alpha}$, где $x^{i}=x^{i}(t)$-уравнения кривой $L$ в нек-рой координатной окрестности ее точки $x_{0}$ с координатами $x^{i}(0)$.



Любые 1-формы $\omega_{0}^{i}, \theta_{i}^{j}, \omega_{i}^{0}$, заданные ва П и удовлетворяющие уравнениям (5) с правыми частями, выражающимися через $\omega_{0}^{k} \wedge \omega_{0}^{l}$, где $\omega_{0}^{i}, i=1, \ldots, n$, линейно независимы, определяют в этом смысле нек-рую П. с. на $M$.



Кривая, к-рую описывает в $\left(P_{n}\right)_{x_{0}}$ точка, определяемая первым вектором $e_{0}(t)$ решения $\left\{e_{\alpha}(t)\right\}$ системы (6), наз. р а з в е р т к о й кривой $L$. Кривая наз. г е о д ез и ч е с к о й ли н и е й П. с. на $M$, если ее развертка в нек-рой окрестности произвольной ее точки $x$ является прямой пространства $\left(P_{n}\right)_{x}$. Уравнения $x^{i=}$ $=x^{i}(t)$ геодезич. линии определяются с помощью функций



из системы



\[

\xi(t)=\left(\omega^{i}\right)_{x(t)}(\dot{x}(t))

\]



\[

d \xi^{i}+\xi^{j}\left(\theta_{j}^{l}\right)_{x(t)}(\dot{x}(t))=\vartheta_{x(t)}(\dot{x}(t)) \xi^{i},

\]



где $\vartheta$ - нек-рая 1-форма. В репере, где $\omega^{i}=d x^{i}$, $\xi^{i}=\dot{x^{i}}$, эта система имеет вид



\[

\begin{gathered}

\frac{d^{2} x^{a}}{\left(d x^{n}\right)^{2}}=-Q^{a}\left(\frac{d x^{1}}{d x^{n}}, \ldots, \frac{d x^{n-1}}{d x^{n}}\right)+ \\

+\frac{d x^{a}}{d x^{n}} Q^{n}\left(\frac{d x^{1}}{d x^{n}}, \ldots, \frac{d x^{n-1}}{d x^{n}}\right)

\end{gathered}

\]



где $Q^{a}$ и $Q^{n}$ - многочлены 2-го порядка, коәфффициенты к-рых есть функции от $x^{1}, \ldots, x^{n}$.



Т е о р е м а $\boldsymbol{K}$ а р т а н а: если на гладком многообразии $M$ задана система кривых, локально определяемая системой дифференциальных уравнений вида (7), то существует одна и только одна нормальная П. с., для к-рой эта система кривых является системой геодезит. линий.



Теория П. с. дает, таким образом, средство для инвариантного исследования систем дифференциальных уравнений специального вида. П. с. полезны также при исследовании геодезических (или проективных) отображений пространств аффинной связности. П. с. сводится к афбинной связности, если на $M$ существуют локальные поля реперов, относительно к-рых $\omega_{i}^{0}=P_{i j} \omega_{0}^{j}$. Для каждой аффинной связности на $M$ существует единственная нормальная П. с., имеющая общие геодезич. линии, из к-рой она может быть получена. Две аф̆ффинные связности геодезически (или проективно) эквйвалентны, если их нормальные П. с. совпадают. В частности, аффинная связность на $M$ при $\operatorname{dim} M>2$ проективно евклидова тогда и только тогда, когда ее тензор проективной кривизны $K_{i k l}$ обрашается в нуль.



Jum.: [1] C a r t a n E., "Bull. Soc. math. France", 1924, t. 52, p. 205-41; 2] e ro жe, Lecons sur la theorie des espaces аффинной, проективной и конформной связности, пер с франц., Казань, 1962; [4] K o b a y a s h i S., N a g a n a T , "J Math. and Mech.", $1964, \mathrm{v} .13$, №, 2, p. 215-35.



ПРОЕКТИВНАЯ СХЕМА - замкнутая подсхема проективного пространства $P_{k}^{n}$; в однородных координатах $x_{0}, \ldots, x_{n}$ на $P_{k}^{n}$ проективная схема задается системой однородных алгебраич. уравнений:



\[

f_{1}\left(x_{0}, \ldots, x_{n}\right)=0, \ldots, f_{r}\left(x_{0}, \ldots, x_{n}\right)=0 .

\]



Каждая П. с. является полной (компактной в случае $k=\mathbb{C}) ;$ обратно, полная схема проективна, если на ней есть обильный обратимый пучок. Имеются и другие критерии проективности.



Обобщением понятия П. с. служит проективный морфизм. Морфизм схем $f: X \rightarrow Y$ наз. п р ое к т и в ны м (а $X$ - схемой, проективной над $Y$ ), если $X$ является замкнутой подсхемой проективного расслоения $P_{Y}(\mathscr{尺})$,



\includegraphics[max width=\textwidth, center]{2023_07_08_f8d7f934cb1071607f8eg-1}

проективных морфизмов іроективна. Проективность морфизма сохраняется и при замене базы; в частности, слои проективного морфизма являются проективными схемами (во не обратно). Если схема $X$ проективна, а $X \rightarrow Z$ - конечный сюръективный морфизм, то и $Z$ проективна.



Любая П. с. (над $Y$ ) может быть получена при помощи конструкции п р о е к т и в н о г о с п е к т р a. Ограничиваясь случаем аффинной базы, $Y=\operatorname{Spec} R$;

пусть $A=\bigoplus_{i \geqslant 0} A_{i}-$ градуированная $R$-алгебра, причем $R$-модуль $A_{1}$ имеет конечный тип и порождает ал-

гебру $A$, и пусть $\operatorname{Proj}(A)-$ множество однородных простых идеалов $p \subset A$, не содержащих $A_{1}$. Снабженное естественной топологией и структурным пучком множество Proј $(A)$ является проективной $Y$-схемой; более того, любая проективная $Y$-схема имеет такой вид.



Лuт.: [1] М а м ф о р д Д., Алгебраическая геометрия, М. , 1979; [2] Итоги науки и техники. Алгебра. Топология. Геометрия, т. 10, М., 1972, с. $47-112$. Алгера. Топология. Гео IРОЕКТИВНОЕ АЛГЕБРАИЧЕСКОЕ МНОЖЕСТВо - подмножество точек проективного пространства $P^{n}$, определенного над полем $k$, пмеющее (в однородвых координатах) вид



$V(I)=\left\{\left(a_{0}, \ldots, a_{n}\right) \in P^{n} \mid \quad f\left(a_{0}, \ldots, a_{n}\right)=0\right.$



для любого $f \in I\}$.



Здесь $I$ - однородный идеал в кольце мвогочленов $k\left[X_{0}, \ldots, X_{n}\right]$. (Идеал $I$ однороден, если из $f \in I$ и $f=$ $=\sum f_{i}$, где все $f_{i}$-однородные многочлены степени $i$, следует, что все $\left.f_{i} \in I.\right)$



Свойства П. а. м.:



\begin{enumerate}

  \item $V\left(\sum_{i \in S} I_{i}\right)=\bigcap_{i \in S} V\left(I_{i}\right)$



  \item $V\left(I_{1} \cap I_{2}\right)=V\left(I_{1}\right) \cup V\left(I_{2}\right)$;



  \item если $I_{1} \subset I_{2}$, то $V\left(I_{2}\right) \subset V\left(I_{1}\right)$;



  \item $V(I)=V(\sqrt{I})$,



\end{enumerate}



где $\sqrt{ } \frac{1}{-}$ радикал идеала $I$. Из свойств 1) -3) следует, что на $V(I)$ можно ввести топологию Зариского. Если $I=\sqrt{I}$, то $I$ однозначно представляется в виде пересечения однородных простых идеалов:



\[

I=\mathfrak{B}_{1} \cap \ldots \cap \mathfrak{B}_{s}

\]



и



\[

V(I)=V\left(\mathfrak{B}_{1}\right) \cup \ldots \cup V\left(\mathfrak{B}_{s}\right) .

\]



В случае, когда I-однородный простой идеал, П. а. м. $V(I)$ наз. п р о е к т и в ны м много обр а з и е м.



Лит.: [1] Іі а ф а р е в и ч И. Р., Основы алгебраической геометрии, М., 1972; [2] 3 а р и с с к и и О., С а м ю э Ј ь П., Коммутативная алгебра, пер. с англ., т. 1-2, M., 1963 .



IІРОЕКТИВНОЕ ИЗГИБАНИЕ - распростране на проективную геометрию понятия изгибания (наложения) в метрич. теории поверхностей, дано Г. Фубини (G. Fubini, 1916) (обобщение этого понятия на геометрию любой группы преобразований получил Э. Картан, E. Cartan, 1920) с использованием понятия т. н. качения одной поверхности по другой.



Пусть $G$ - группа преобразований пространства $E$. Поверхность $S^{\prime}$ налагается на поверхность $S$ (или катится по $S$ ) в геометрии группы $G$, если между их точками устанавливается взаимно однозначное соответствие так, что каждой паре соответствующих точек $M \in S$, $M^{\prime} \in S^{\prime}$ можно присоединить преобразование $y \in G$, к-рое переводит $S^{\prime}$ в положение $S$. При этом требуется, чтобы



\begin{enumerate}

  \item $M^{\prime}$ совмещалась с $M$;



  \item каждая кривая $l^{\prime} \in S^{\prime}$, проходящая через $M$, имела в этой точке касавие $n$-го порядка с соответствующей



\end{enumerate}



\end{document}